\documentclass{../../../templates/natsirt}

\usepackage{pgfplots}
\pgfplotsset{compat=1.18}
\tcbuselibrary{breakable}
\tcbset{breakable}

\title{Probability}
\author{Yi Hong}
\subtitle{Note:p}

\setlength{\parindent}{0pt}
\setlength{\parskip}{3pt}
\setlength{\headheight}{15pt}

% Custom Box Environments for Pure Active-Recall Checklist
\newcounter{topic}[section]

\newenvironment{recallbox}{%
    \begin{bluebox}\refstepcounter{topic}\textbf{Definitions \& Concepts to State \thesection.\thetopic}%
}{\end{bluebox}}

\newenvironment{theoremsbox}{%
    \begin{greenbox}\refstepcounter{topic}\textbf{Theorems \& Identities to State \thesection.\thetopic}%
}{\end{greenbox}}

\newenvironment{challengebox}{%
    \begin{redbox}\refstepcounter{topic}\textbf{Proof Challenges \& Calculations \thesection.\thetopic}%
}{\end{redbox}}

\newenvironment{warningbox}{%
    \begin{whitebox}\textbf{Conventions \& Exam Pitfalls. }%
}{\end{whitebox}}

\newcommand{\citem}{\item[$\square$]}
\newcommand{\citemdone}{\item[$\text{\rlap{$\checkmark$}}\square$]}

\DeclareMathOperator{\Var}{Var}
\DeclareMathOperator{\Bin}{Bin}
\DeclareMathOperator{\Poi}{Poi}
\DeclareMathOperator{\Geo}{Geo}
\DeclareMathOperator{\NegBin}{NegBin}

\begin{document}

% =========================================================================
% CHAPTER 1: BASIC COMBINATORICS
% =========================================================================
\section{Basic Combinatorics}

\begin{warningbox}
\begin{itemize}
    \citem This chapter discusses combinatorics topics that fully overlap with the chapter taught in Discrete Mathematics; refer back to the Discrete Mathematics notes.
\end{itemize}
\end{warningbox}


% =========================================================================
% CHAPTER 2: AXIOMS OF PROBABILITY
% =========================================================================
\section{Axioms of Probability}

\begin{warningbox}
\begin{itemize}
    \citem When a question says ``of these balls are withdrawn one at a time'' (or similar), it usually involves permutations.
\end{itemize}
\end{warningbox}

\begin{recallbox}
\begin{itemize}
    \citem Write the definition of sample space.
    \citem Write the definition of an event.
    \citem Write the definition of the union and intersection of events.
    \citem Write the definition of mutually exclusive events.
    \citem Write the definition of the null event.
    \citem Write the definition of the complement of an event.
    \citem Write the definition of the difference of two events (sets).
    \citem What does it mean for event $E$ to be contained in $F$ ($E \subseteq F$)?
\end{itemize}
\end{recallbox}

\begin{theoremsbox}
\begin{itemize}
    \citem Write De Morgan's Laws.
    \citem Write the three Axioms of Probability.
    \citem Write the probability version of the Principle of Inclusion-Exclusion (PIE).
\end{itemize}
\end{theoremsbox}

\begin{challengebox}
\begin{itemize}
    \citem Prove that $P(\emptyset) = 0$.
    \citem Prove that Axiom 3 also works for the finite case.
    \citem In the experiment of tossing a coin, if the coin is biased and we believe that a head is twice as likely to appear as a tail, find $P(\{H\})$ and $P(\{T\})$. (Use the axioms to do it.)
    \citem Prove the Complement Rule:
    \[
        P(E^c) = 1 - P(E).
    \]
    \citem Prove that if $E \subseteq F$, then $P(E) \le P(F)$.
    \citem Prove the Addition Rule: if $E$ and $F$ are two events of an experiment,
    \[
        P(E \cup F) = P(E) + P(F) - P(E \cap F).
    \]
    \citem Prove
    \[
        P(B \setminus A) = P(B) - P(B \cap A).
    \]
    \citem Assume that $S$ is the sample space of an experiment with $|S|$ equally likely outcomes. Prove that if $E$ is an event of the experiment, then
    \[
        P(E) = \frac{|E|}{|S|}.
    \]
    \citem A poker hand consists of 5 cards. A full house is a three of a kind plus a pair. What is the probability that one is dealt a full house?
    \citem In the game of bridge, the entire deck of 52 cards is dealt out to 4 players Amy, Betsy, Christine, and Diana. Find the probability that one of them receives all 13 spades.
    \citem In the game of bridge, the entire deck of 52 cards is dealt out to 4 players $A, B, C, D$. Find the probability that each player receives 1 ace.
\end{itemize}
\end{challengebox}


% =========================================================================
% CHAPTER 3: CONDITIONAL PROBABILITY AND INDEPENDENCE
% =========================================================================
\section{Conditional Probability and Independence}

\begin{recallbox}
\begin{itemize}
    \citem Write the definition of conditional probability.
    \citem Write the definition of independent events for two events $E$ and $F$ ($P(E \mid F) = P(E)$).
    \citem Write the definition of independence for three events $E, F, G$ (state all four required conditions).
    \citem Write the definition of independence for $n$ events $E_1, E_2, \dots, E_n$.
\end{itemize}
\end{recallbox}

\begin{theoremsbox}
\begin{itemize}
    \citem Notice the Generalization of the Multiplication Rule: if $E_1, E_2, \dots, E_n$ are events of an experiment,
    \[
        P(E_1 E_2 \cdots E_n) = P(E_1)\, P(E_2 \mid E_1)\, P(E_3 \mid E_1 E_2) \cdots P(E_n \mid E_1 E_2 \cdots E_{n-1}).
    \]
    \citem Write Bayes' Formula.
    \citem Notice the Generalization of the Law of Total Probability: if $F_1, \dots, F_n$ are mutually exclusive events such that $\bigcup_{k=1}^n F_k = S$, then
    \[
        P(E) = \sum_{k=1}^n P(E \cap F_k) = \sum_{k=1}^n P(E \mid F_k) P(F_k).
    \]
    \citem Based on the above, write the Generalization of Bayes' Formula.
    \citem Notice the symmetry of independence: if $E$ is independent of $F$, then $F$ is independent of $E$.
\end{itemize}
\end{theoremsbox}

\begin{challengebox}
\begin{itemize}
    \citem In the card game bridge, the 52 cards are dealt out equally to 4 players --- Alice, Bob, Claire, and Danny. If Alice and Bob have a total of 8 spades among them, what is the probability that Claire has 3 of the remaining spades?
    \citem Prove that if $E$ and $F$ are events of an experiment,
    \[
        P(E \cap F) = P(E \mid F) P(F).
    \]
    \citem Prove that if $E$ and $F$ are events of an experiment, then
    \[
        P(E) = P(E \mid F) P(F) + P(E \mid F^c) P(F^c).
    \]
    \citem At a certain stage of a criminal investigation, the inspector in charge is 60\% convinced of the guilt of a certain suspect. However, a new piece of evidence shows that the criminal has a certain characteristic. If 20\% of the population possesses this characteristic, how certain of the guilt of the suspect should the inspector now be if it turns out that the suspect has the characteristic?
    \citem Twins can be identical or fraternal. If identical twins must be of the same sex, and fraternal twins have a 50\%--50\% chance of having the same sex or not, and some data indicated that 64\% of twin births are same-sexed, what is the percentage of identical twins?
    \citem Prove that two events $E$ and $F$ are independent if and only if
    \[
        P(E \cap F) = P(E) P(F).
    \]
    \citem Prove that if $E$ and $F$ are independent events, then $E$ and $F^c$ are also independent events.
    \citem Prove that if $E$, $F$, and $G$ are independent events, then $E$ and $F \cap G$ are independent events.
    \citem Two fair dice are thrown. Let $E$ denote the event that the sum of the dice is 7, $F$ denote the event that the first die equals 4, and $G$ denote the event that the second die equals 3. Are the three events $E$, $F$, and $G$ independent?
    \citem Independent trials consisting of rolling a pair of fair dice are performed. The outcome of a roll is the sum of the dice. What is the probability that an outcome of 5 appears before an outcome of 7?
\end{itemize}
\end{challengebox}


% =========================================================================
% CHAPTER 4: RANDOM VARIABLES
% =========================================================================
\section{Random Variables}

% -------------------------------------------------------------------------
% SECTION 4.1: RANDOM VARIABLES
% -------------------------------------------------------------------------
\subsection{Random Variables}

\begin{recallbox}
\begin{itemize}
    \citem Write the definition of a random variable $X$.
    \citem How do we define $P(X = x)$?
\end{itemize}
\end{recallbox}

\begin{challengebox}
\begin{itemize}
    \citem A life insurance agent has two elderly clients, $A$ and $B$, each of whom has a life insurance policy that pays 1 million ringgit upon death. Let $E_A$ be the event that $A$ will die in the following year, and $E_B$ be the event that $B$ will die in the following year. Assume that $E_A$ and $E_B$ are independent, with probabilities $P(E_A) = 0.05$ and $P(E_B) = 0.1$. Let $X$ be the total amount of money (in units of 1 million ringgit) that will be paid out this year to any of these clients' beneficiaries. Find the probability distribution of $X$.
\end{itemize}
\end{challengebox}

% -------------------------------------------------------------------------
% SECTION 4.2: DISCRETE RANDOM VARIABLES
% -------------------------------------------------------------------------
\subsection{Discrete Random Variables}

\begin{recallbox}
\begin{itemize}
    \citem Write the definition of a discrete random variable.
    \citem Write the definition of the probability mass function (PMF).
    \citem Write the definition of the cumulative distribution function (CDF).
    \begin{itemize}
        \item Remark: The domain of the function is always the real numbers ($\R$), even though the random variable is discrete.
    \end{itemize}
\end{itemize}
\end{recallbox}

\begin{theoremsbox}
\begin{itemize}
    \citem Notice that for the PMF $p(x)$, it satisfies:
    \begin{itemize}
        \item $0 \le p(x) \le 1$
        \item $\sum_x p(x) = 1$
    \end{itemize}
\end{itemize}
\end{theoremsbox}

% -------------------------------------------------------------------------
% SECTION 4.3: EXPECTED VALUE
% -------------------------------------------------------------------------
\subsection{Expected Value}

\begin{recallbox}
\begin{itemize}
    \citem Write the definition of the expected value of a discrete random variable.
\end{itemize}
\end{recallbox}

% -------------------------------------------------------------------------
% SECTION 4.4: EXPECTATION OF A FUNCTION OF A RANDOM VARIABLE
% -------------------------------------------------------------------------
\subsection{Expectation of a Function of a Random Variable}

\begin{recallbox}
\begin{itemize}
    \citem Write the definition of the $n$th moment of $X$.
\end{itemize}
\end{recallbox}

\begin{challengebox}
\begin{itemize}
    \citem Prove that if $X$ is a discrete random variable with probability mass function $p(x)$, then for any real-valued function $g$ of a real variable,
    \[
        E(g(X)) = \sum_x g(x) p(x).
    \]
    \citem Prove that if $a$ and $b$ are constants, then
    \[
        E(aX + b) = aE(X) + b.
    \]
    More generally, for any two real-valued functions $g_1$ and $g_2$,
    \[
        E(g_1(X) + g_2(X)) = E(g_1(X)) + E(g_2(X)).
    \]
\end{itemize}
\end{challengebox}

% -------------------------------------------------------------------------
% SECTION 4.5: VARIANCE
% -------------------------------------------------------------------------
\subsection{Variance}

\begin{recallbox}
\begin{itemize}
    \citem Write the definition of the variance of $X$ and standard deviation.
\end{itemize}
\end{recallbox}

\begin{challengebox}
\begin{itemize}
    \citem Prove that if $X$ is a discrete random variable,
    \[
        \Var(X) = E(X^2) - (E(X))^2.
    \]
    \citem Prove that if $a$ and $b$ are constants, then
    \[
        \Var(aX + b) = a^2 \Var(X).
    \]
\end{itemize}
\end{challengebox}

% -------------------------------------------------------------------------
% SECTION 4.6: BERNOULLI AND BINOMIAL RANDOM VARIABLES
% -------------------------------------------------------------------------
\subsection{Bernoulli and Binomial Random Variables}

\begin{recallbox}
\begin{itemize}
    \citem Write the definition of a Bernoulli random variable.
    \citem Write the definition of a binomial random variable and its corresponding PMF. How do we denote that $X$ is a binomial random variable with parameters $n$ and $p$?
\end{itemize}
\end{recallbox}

\begin{challengebox}
\begin{itemize}
    \citem If $X$ is a binomial random variable with parameters $n$ and $p$, prove that
    \begin{gather*}
        E(X) = np, \\
        \Var(X) = np(1 - p).
    \end{gather*}
    \citem Prove that if $X \sim \Bin(n, p)$, where $0 < p < 1$, then as $k$ goes from $0$ to $n$, $P(X = k)$ first increases monotonically and then decreases monotonically, reaching its largest value when $k$ is the largest integer less than or equal to $(n + 1)p$.
\end{itemize}
\end{challengebox}

% -------------------------------------------------------------------------
% SECTION 4.7: POISSON RANDOM VARIABLE
% -------------------------------------------------------------------------
\subsection{Poisson Random Variable}

\begin{recallbox}
\begin{itemize}
    \citem Write the definition of a Poisson distribution and its notation.
\end{itemize}
\end{recallbox}

\begin{challengebox}
\begin{itemize}
    \citem Prove that if $X \sim \Bin(n, p)$ with $n$ large and $p$ small (rare events), then $X$ can be approximated by $\tilde{X} \sim \Poi(\lambda)$, where $\lambda = np$.
    \citem Prove that if $X \sim \Poi(\lambda)$, then
    \begin{gather*}
        E(X) = \lambda, \\
        \Var(X) = \lambda.
    \end{gather*}
    \citem Suppose that the number of typographical errors on a single page of this book has a Poisson distribution with an average of $0.5$ errors per page.
    \begin{itemize}
        \item[(a)] Find the probability that there is at least one error on a page.
        \item[(b)] Find the probability that there is at most one error on a page.
    \end{itemize}
    \citem In a certain region, on average there occur 2 earthquakes in a week. Suppose that the number of earthquakes that occur is in accordance with a Poisson process.
    \begin{itemize}
        \item[(a)] What is the average number of earthquakes in two weeks' time?
        \item[(b)] Find the probability that at least 3 earthquakes occur during the next two weeks.
        \item[(c)] Let $T$ be the time (in weeks), starting from now, until the next earthquake. Find $P(T \le t)$, where $t > 0$.
    \end{itemize}
\end{itemize}
\end{challengebox}

% -------------------------------------------------------------------------
% SECTION 4.8: OTHER DISCRETE PROBABILITY DISTRIBUTIONS
% -------------------------------------------------------------------------
\subsection{Other Discrete Probability Distributions}

\begin{recallbox}
\begin{itemize}
    \citem Write the definition of a geometric random variable.
    \citem Write the definition of a negative binomial random variable.
\end{itemize}
\end{recallbox}

\begin{challengebox}
\begin{itemize}
    \citem If $X \sim \Geo(p)$, prove that
    \begin{gather*}
        E(X) = \frac{1}{p}, \\
        \Var(X) = \frac{1 - p}{p^2}.
    \end{gather*}
    \citem If $X \sim \NegBin(r, p)$, prove that
    \begin{gather*}
        E(X) = \frac{r}{p}, \\
        \Var(X) = \frac{r(1 - p)}{p^2}.
    \end{gather*}
    \citem Find the expected value and the variance of the number of times one must throw a fair die until the outcome 1 has occurred 4 times.
    \citem Independent Bernoulli trials resulting in a success with probability $p$ and a failure with probability $1 - p$ are performed.
    \begin{itemize}
        \item[(a)] What is the probability that 2 successes occur before 3 failures?
        \item[(b)] What is the probability that 3 successes occur before 2 failures?
        \item[(c)] What is the probability that $n$ successes occur before $m$ failures?
    \end{itemize}
\end{itemize}
\end{challengebox}

\end{document}
